




%%%%%%%%%%%%%%%%%%%%%%%%%%%%%%%%%%%%%%%%%%%%%%%%%%%%%%%%%%
%%%
%%%  Completely Simple Semigroups, Lie Algebras, 
%%%   and the Road Coloring Problem 
%%%	
%%%%%%%%%%%%%%%%%%%%%%%%%%%%%%%%%%%%%%%%%%%%%%%%%%%%%%%%%%


\section{Completely Simple Semigroups, Lie Algebras, 
  and the Road Coloring Problem }
\label{sec:4.8}

Studying the dynamics inherent in state-transition diagrams, either in the stochastic setting as Markov chains or in deterministic models as automata, continues to be a rich source of interesting mathematics and applications. 
Even for finite state systems, fascinating unsolved problems exist.
\^Cern\'y's conjecture concerning the length of a synchronizing instruction for deterministic automata is one such example. 
Another open problem in this area is referred to as the ``road coloring problem'' and was the initial motivation for this work. 
The road coloring problem appeared in published form for the first time in a paper by Adler, Goodwyn, and Weiss \cite{AGW} and can be stated very simply:  Suppose you are given a strongly connected directed graph (or state-transition diagram).  
Is there a way to label the edges so that a synchronizing  instruction exists? 
A precise definition of synchronizing instruction will be provided below. 
In \cite{AGW},  it was shown that aperiodicity was necessary for such an instruction to exist. 
The road coloring problem conjectures that aperiodicity is a sufficient condition as well. 

\subsection{Background} 
\label{sec:4.8.1}

A finite automaton consists of a finite set $V$ of states, a finite set of inputs $A$, and a state transition function $\dl :V\times A \to V$.
We can associate with each $a \in A$ a corresponding transformation $R_a: V \to V$ defined by $R_a(v) = \dl(v,a)$, which we write compactly as $va$.
The semigroup $S$ generated under composition by $\{R_a: a \in A\}$ is called the semigroup of the automaton.
Given any automaton, one obtains a labelled digraph by taking the vertices of the graph to be the set of states $V$ and for each $(v,a) \in V \times A$ defining a directed edge from $v$ to $va$ with label $a$.  
Note that the constructed directed graph has \textit{uniform outdegree}, i.e., each vertex of $V$ has the same number of edges leading from it.

If one considers an unlabelled digraph $G = (V,\E)$ of uniform outdegree $d$, then any labelling of the edges with members of $A$, where $|A| = d$, such that each edge issuing from any given vertex has a distinct label is called a \textit{coloring} of the digraph.
Any such coloring uniquely determines an automaton $\dl: V \times A \to V$, where $R_a(v) = va$ is the terminal point of the directed edge with initial point $v$ and label $a$.  
We henceforth assume that digraphs have uniform outdegree, identifying the coloring with the set $\C = \{R_a: a \in A\}$, and we refer to the semigroup $S = \lr \C \rlr$, generated by $\C$, as the \textit{coloring semigroup} for the given labelling.

\begin{definition}   \label{defn:4.6}
Let $S$ be a coloring semigroup of a digraph $G = (V,\E)$.
The \textit{rank} of $s \in S$ is the cardinality $|Vs|$ of the image of $s$.
Any finite transformation semigroup $S$ has a minimal ideal or \textit{kernel}, which consists of the elments of minimal rank (see \cite{H}).
This common minimal rank is called the \textit{rank of the kernel}.
\end{definition}

A \textit{synchronizing instruction} for an edge-labelled digraph is thus any transformation of rank one.
The road coloring problem can now be stated in the following way.

\runinhead{Road Coloring Conjecture.}
\textit{Let $G = (V,\E)$ be a $d$-out, strongly connected digraph.
Let $\{S_i\}$ be the set of all coloring semigroups.  
If $G$ is aperiodic, then for some $i$, $S_i$ contains a synchronizing instruction.} \\

The notion of period used here is the standard one from Markov chain theory, where for any vertex $v$ in a directed graph, $\hbox{per}(v) - \gcd\{|c_1|, \ldots , |c_n|\}$ where $c_1, \ldots , c_n$ are the simple cycles containing $v$.
The assumption of strong connectivity ensures that all vertices have the same period.

\runinhead{Notation.}  We use rank and trace to denote the rank and trace respectively of a matrix.
For a function, rank and trace refer to its representing matrix.

\subsection{Semigroups and kernels}
\label{sec:4.8.2}

Let $S$ be a coloring semigroup of a strongly connected digraph with kernel $K$.
Let $E(S)$ be the set of idempotents of $S$. 
Impose on $E(S)$ the order $e \le f$ iff $ef = fe = e$.
Form the Rees product $x \times G \times Y$ of $K$ in the usual manner choosing an idempotent, $e_0$, minimal with respect to this order and letting
\[
  x = E(Ke_0), \quad G = e_0Ke_0, \quad Y = E(e_0K)
\]
with product
\[
  (x_1,g_1,y_1)(x_2,g_2,y_2) := (x_1,g_1(y_1x_2)g_2,y_2)\,.
\]
Notice that $y_1x_2 \in G$.
The sandwich function
\[
  \phi: Y \times X \to G, \qquad \phi(y,x) = yx
\]
plays an important r\^ole in this paper.
In particular, recall that if $\phi(y,x) = e_0$ for all $(y,x) \in Y \times X$, then $X \times G \times Y$ is a direct product.

Recall that if $K$ is the kernel of a finite transformation semigroup $S$, the sets $X$, $G$, and $Y$ can be further characterized.
Suppose $S$ is a semigroup of functions on the set $\{1, \ldots , n\}$.
Then each idempotent $e \in X$ can be associated with a partition $\pi_e$ of the set $\{1, \ldots , n\}$ in the standard way: $\pi_e = \{e^{-1}(j): 1 \le j \le n \}$.
Thus the set $X$ can also be thought of as a collection of partitions of the set  $\{1, \ldots , n\}$.
Similarly each idempotent $f \in Y$ can be associated with its range $R_f = \{jf: 1 \le j \le n\}$, where we employ the convention of writing inputs to the transformation on the left.
The group $G = e_0Ke_0$ can then be seen as the subgroup of transformations within $K$ having the same partition and range as $e_0$\,.

The ideal structure of $K$ will also play a r\^ole in our discussions.
Each set of the form $eK$ with $e \in X$ is a minimal right ideal consisting of the set of all transformations in $K$ having partition $\pi_e$\,.
Similarly, $Kf$ with $f \in Y$ is a minimal left ideal consisting of the set of all transformations in $K$ having range $R_f$\,.
When $X = E(Ke_0) = \{e_0\}$, a single idempotent, $K$ is, of course, a right group.
Notice that when a coloring semigroup has a synchronizing instruction or rank one instruction, $K$ is easily shown to be a right group.

The importance of right groups in the context of the road colroing problem has been shown in recent papers \cite{B, CKK} and this paper continues an exploration of their r\^ole.
In fact, we have the following theorem.

%
\begin{theorem}    \label{thm:4.13}
Suppose a strongly connected, apreiodic digraph $G$ has a coloring semigroup $S$ whose kernel $K$ is a right group.
Then $G$ has a coloring semigroup that contains a synchronizing instruction.
\end{theorem}
%

%
\begin{proof}   
Suppose the (unique) vertex partition induced by the right group $K$, is $\Pi = \{P_1,P_2 \ldots P_l\}$, where rank$(K) = t \ge 2$ (as we are done if t=1).

Notice that every element of $S$, and therefore the generators of $S$, act as permutations on the elements of $\Pi$.
Suppose not.
Then there are vertices $v_i \in P_i$, $v_j \in P_j$, and an element $s$ in $S$ such that
\[
  v_is \in P_m, \qquad v_js \in P_m\,.
\]
Then, if $k$ is any element of $K$,
\[
  v_isk \in P_mk, \qquad v_jsk \in P_mk\,.
\]
But $P_mk = w$, since the elements of the kernel collapse the partition elements to a single vertex.
Since $sk$ is in the kernel, the $\{v: vsk = w\}$ is one of the partition elements.  
Thus $P_i = P_j$, showing that $s$ is a permutation on partition elements.

Form the quotient graph induced by the coloring on the ``vertices'' $\Pi = \{P_1,P_2, \ldots  P_l\}$.
Since the generators act as permutations on the elements of $\Pi$, the resulting graph will be ``Eulerian'', having the indegree of each vertex equal to its outdegree.
Now, we will use the algorithm introduced by Kari \cite{K}, to find a coloring that synchronizes this Eulerian graph.
This is done in the following way.

Choose any element $w_0 = c_1c_2 \cdots c_l$  from the original kernel $K$, each $c_i$ being a generator of $S$.
Note that $P_1w_0$ is a single vertex in the original graph.
Proceeding as in \cite{K}, choose a vertex $P_{i_0}$, in the quotient graph, such that for colors $r$ and $b$,
\[
  P_{i_0}r \not= P_{i_0}b\,.
\]

For step one, swap the labels of the $r$-edges leaving $P_{i_0}$ with the $b$-edges leaving $P_{i_0}$.
No other edge labels are changed.

For step two, we reassign colors in $w_0$ as follows.
If $P_1 = P_{i_0}$ then change $c_1$.
If $c_1$ is either $r$ or $b$, swap $r$ with $b$, otherwise leave it unchanged.
Working through $w_0$, whenever $P_1c_1c_2 \cdots c_{j-1} \subset P_{i_0}$, change $c_j$ accordingly.
This results in a new word, $w_1 = c'_1c'_2 \cdots c'_l$, with $P_1w1$ remaining a single vertex.

Each stage of this procedure is performed in the same way.
At stage $k$, Kari's algorithm guarantees that the new quotient graph is Eulerian with fewer vertices.
Step one consists of finding a vertex in this new quotient graph and relabelling as described above.
For step two, the word $w_{k-1}$ is changed each time the vertices of $P_1$ map to that partition element that was changed by this $k$\,th relabelling, ensuring that $P_1w_k$ is a single vertex.

This procedure terminates in a synchronizing instruction $w'$ for the quotient graph induced by the kernel, i.e., $w'$ maps every $P_j$ to a single $P_{j_0}$.
Since the graph is strongly connected at the quotient level, there exists an instruction $w''$ such that $P_{j_0}w'' \subset P_1$.
If the procedure requires $m$ steps then $w_{\hbox{\small sync}} = w'w''w_m$ is a synchronizing instruction for the original graph.
\qed
\end{proof}
